\subsection{局部整 Noether 环；正规 Noether 环}

设 A 为 Noether 环。用 \((\mathrm{Y}_{i})_{i\in\mathrm{I}}\) 表示 Spec(A) 的连通分支组成的有限族 (II, § 4, 第 2 节, 命题 10 及第 3 节, 命题 11 的推论 7)。根据 II, § 4, 第 3 节, 命题 15，对每个 i 存在 A 的唯一幂等元 \(e_{i}\) 使 \(\mathrm{Y}_{i}=\mathrm{V}(e_{i})\) ，且从 A 到诸环 A/Ae \(_{i}\) 之积的典范映射是双射。A 的商环 A/Ae \(_{i}\) 称为 A 的典范分量。令 \(f_{i}=1-e_{i}\) 。有 \(\sum_{i\in I}f_{i}=1\) ，且 \((f_{i})_{i\in I}\) 是 A 的非零幂等元组成的正交族（同上引文）。由此得出，\(f_{i}\) 在 A/Ae \(_{j}\) 中的像当 j=i 时等于 1，当 j≠i 时等于 0；于是由环同态 A→ \(\prod_{j}A/Ae_{j}\) 得出典范同构 A \(_{f_{i}}\rightarrow A/Ae_{i}\) 。

根据同上引文, 命题 14 的推论 2，下列条件等价： Spec(A) 的连通分支是不可约的；

\begin{enumerate}
\def\labelenumi{(\roman{enumi})}
\setcounter{enumi}{1}
\item
  A 的每个素理想（相应地，极大理想）只属于 Spec(A) 的一个不可约分支；
\item
  A 的每个素理想（相应地，极大理想）只包含一个极小素理想；
\item
  对 A 的每个素理想（相应地，极大理想）p，拓扑空间 \(\operatorname{Spec}(A_{\mathfrak{p}})\) 是不可约的；
\item
  对 A 的每个典范分量 C，拓扑空间 \(\operatorname{Spec}(\mathbb{C})\) 是不可约的。
\end{enumerate}

现在注意，若 A 是约化的，则所有环 \(A_{p}\) 都是约化的 (II, § 2, 第 6 节, 命题 17)；反之，若对 A 的每个极大理想 m，环 \(A_{m}\) 都约化，则 A 约化 (II, § 3, 第 3 节, 定理 1 的推论 2)。于是应用 II, § 4, 第 3 节, 命题 14 的推论 1，由前面所述推出下列条件的等价性：

\begin{enumerate}
\def\labelenumi{(\roman{enumi})}
\item
  A 约化且 Spec(A) 的连通分支不可约；
\item
  对 A 的每个素理想（相应地，极大理想）p，环 A \(_{p}\) 是整环；
\item
  A 的典范分量都是整环。
\end{enumerate}

若一个 Noether 环满足上述等价条件 至 ，就称它是局部整的。

设环 A 局部整；设 u 为 A 到整环的（有限）积 \(\prod_{j\in J}A_{j}\) 上的同构。存在双射 \(\sigma:J\to I\) ，使与 u 相伴的从 \(\operatorname{Spec}(\prod_{j\in J}A_{j})\) 到 \(\operatorname{Spec}(A)\) 的映射给出 \(\operatorname{Spec}(A_{j})\) 到连通分支 \(Y_{\sigma(j)}\) —— \(\operatorname{Spec}(A)\) 的连通分支——上的同胚；于是由 u 得出典范分量 \(A/Ae_{\sigma(j)}\) 到 \(A_{j}\) 上的同构。

命题 14. - 设 A 为 Noether 环。下列条件等价：

\begin{enumerate}
\def\labelenumi{(\roman{enumi})}
\item
  A 约化且在其全分式环中整闭；
\item
  A 同构于一个有限族整闭环的积。
\item
  A 的典范分量是整闭的；
\item
  对 A 的每个素理想（相应地，极大理想）p，环 \(A_{p}\) 整闭。
\end{enumerate}

 与 的等价性由 V, § 1, 第 2 节, 命题 9 的推论 2 得出， 与 的等价性由命题前的注记得出。设 p 为 A 的素理想；存在 A 的唯一典范分量 A' 使 p 属于 Spec(A) 的闭子集 Spec(A')，且我们有典范同构 \(A_{p} \simeq A'_{p}\)。于是 与 的等价性由同上引文, 第 5 节, 命题 16 的推论 1 与推论 3 得出。

若环 A 是 Noether 的且满足命题 14 的等价条件 至 ，则称 A 为正规环。Noether 环整闭当且仅当它是整环且正规。正规局部环是整闭的。

